\subsection{紧扰动与本质谱}

本条中，E 是一个无限维复希尔伯特空间。

引理 13. --- 设 I 是 E 中的标准正交族。沿 I 的有限子集补集滤子，族 I 弱收敛于 0。

设 \(x \in E\)。根据 Bessel 不等式（EVT，第 V 卷，第 21 页，命题 4）和 TG，第 IV 卷，第 37 页定理 1，族 \((|\langle e_{i} | x\rangle|^{2})_{i\in I}\) 可求和，故 \(\langle e_{i} | x\rangle\) 沿 I 的有限子集补集滤子趋于 0（TG，第 III 卷，第 38 页，命题 1）。

命题 18. --- 设 \(u\) 是 E 上的自伴部分算子，\(\lambda\) 是实数。有 \(\lambda \in \mathrm{Sp}_{\mathrm{e}}(u)\) 当且仅当 E 中存在一个标准正交序列 \((x_{n})_{n \in \mathbf{N}}\)，使得 \(u(x_{n}) - \lambda x_{n}\) 在 E 中趋于 0。

首先假设 \(\lambda \in \mathrm{Sp}_{\mathrm{e}}(u)\)。若 \(u\) 关于 \(\{\lambda\}\) 的谱投影具有无限秩，则其像中的每个标准正交序列 \((x_{n})\) 都满足要求（参见 IV 第 278 页命题 9 的推论）。以下假设情形并非如此。

对每个 \(k \in \mathbf{N}\)，令 \(\mathrm{J}_k\) 为满足下式的 \(t \in [\lambda - 1, \lambda + 1]\) 的集合：

\[
{\frac {1}{k + 2}} <   | t - \lambda | \leqslant {\frac {1}{k + 1}}.
\]

集合 \(J_{k}\) 两两不交。此外，对每个整数 \(K \in N\)，集合族 \((\mathrm{J}_{k})_{k \geqslant K}\) 构成集合

\[
\mathrm{I} _ {\mathrm{K}} = [ \lambda - 1 / (\mathrm{K} + 1), \lambda + 1 / (\mathrm{K} + 1) ] - \{0 \}.
\]

 由于算子 \(u\) 关于 \(\mathrm{I_K} \cup \{0\}\) 的谱投影具有无限秩（IV 第 297 页引理 9），而 u 关于 \(\{\lambda\}\) 的谱投影按假设具有有限秩，根据 IV 第 278 页命题 8，存在严格递增的序列 \((k_n)_{n \in \mathbf{N}}\)（在 \(\mathbf{N}\) 中），使得谱投影 \(p_{n}\)（u 关于 \(J_{k_{n}}\) 的）非零。令 \(x_{n}\) 为 \(p_{n}\) 的像中范数为 1 的向量。序列 \((x_{n})\) 是标准正交的，因为 \(p_{n}\) 的像与 \(p_{m}\) 的像对于 N 中所有 \(n \neq m\) 都正交。

设 \(n \in N\)。记 \(\mu_{n}\) 为 \(x_{n}\) 关于 u 的谱测度；其支撑包含于 \(J_{k_{n}}\)（IV 第 278 页命题 9）。由此

\[
\| u (x _ {n}) - \lambda x _ {n} \| ^ {2} = \int_ {\mathbf {C}} | t - \lambda | ^ {2} d \mu_ {n} (t) \leqslant \frac {1}{k _ {n} ^ {2}},
\]

因此序列 \((x_{n})_{n\in\mathbf{N}}\) 具有所要求的性质。

反之，假设 E 中存在标准正交序列 \((x_{n})_{n\in \mathbf{N}}\)，使得 \(u(x_{n}) - \lambda x_{n}\to 0\)。记 \(\mu_{n}\) 为 \(x_{n}\) 关于 \(u\) 的谱测度。

令 \(\varepsilon > 0\)。记 \(p_{\varepsilon}\) 为 \(u\) 关于开区间 \(\mathrm{I}_{\varepsilon} = ]\lambda - \varepsilon, \lambda + \varepsilon[\) 的谱投影。对每个 \(n \in \mathbf{N}\)，有

\[
\begin{array}{r l} & 1 = \| x _ {n} \| ^ {2} = \mu_ {n} (\mathrm{I} _ {\varepsilon}) + \mu_ {n} (\mathbf {R} - \mathrm{I} _ {\varepsilon}) \\ & \qquad \leqslant \mu_ {n} (\mathrm{I} _ {\varepsilon}) + \frac {1}{\varepsilon^ {2}} \int_ {\mathbf {R} - \mathrm{I} _ {\varepsilon}} | t - \lambda | ^ {2} d \mu_ {n} (t) \\ & \qquad = \| p _ {\varepsilon} (x _ {n}) \| ^ {2} + \frac {1}{\varepsilon^ {2}} \| u (x _ {n}) - \lambda x _ {n} \| ^ {2}. \end{array}
\]

因此，关于序列 \((x_{n})\) 的假设说明 \(p_{\varepsilon}(x_{n})\) 的范数不可能趋于 0。由于标准正交序列 \((x_{n})\) 在 E 中弱收敛于 0（引理 13），投影 \(p_{\varepsilon}\) 不可能是紧的（III 第 6 页命题 6 的推论），因而具有无限秩。由于这对每个 \(\varepsilon > 0\) 都成立，IV 第 297 页引理 9 使我们可以得出 \(\lambda \in \mathrm{Sp}_{\mathrm{e}}(u)\)。

下面的定理是 III 第 93 页定理 3 对自伴部分算子的对应结果。

定理 4. --- 设 u 是 E 上的自伴部分算子。u 的本质谱是集合 Sp(u + v) 的交，其中 v 遍历 E 的紧埃尔米特自同态的集合。

对 E 的每个埃尔米特自同态 v，部分算子 \(u + v\) 都是自伴的，因为 \((u + v)^{*} = u^{*} + v^{*}\)（参见 IV，第 236 页）。

证明若 \(v\) 是紧的，则 \(\mathrm{Sp_e}(u) \subset \mathrm{Sp_e}(u + v)\)。设 \(\lambda \in \mathrm{Sp_e}(u)\)，并令 \((x_n)_{n \in \mathbf{N}}\) 为 E 中满足 \(u(x_n) - \lambda x_n\) 趋于 0 的标准正交序列（命题 18）。序列 \((x_n)\) 在 E 中弱收敛于 0（引理 13）；由于自同态 \(v\) 是紧的，序列 \((v(x_n))_{n \in \mathbf{N}}\) 在 E 中趋于 0（III 第 6 页命题 6 的推论）。因此，序列 \((u+v)(x_{n})-\lambda x_{n}\) 在 E 中趋于 0，命题 18 蕴含 \(\lambda\in\mathrm{Sp}_{\mathrm{e}}(u+v)\)。

因此，算子 \(u\) 的本质谱包含于 \(\mathrm{Sp}(u + v)\) 的交，其中 \(v \in \mathcal{L}^{\mathrm{c}}(\mathrm{E})\) 遍历埃尔米特元素。

反之，设 \(\lambda \in \mathrm{Sp}_{\mathrm{s}}(u)\)。记 \(\mathrm{E}_{\lambda}\) 为算子 \(u\) 关于 \(\lambda\) 的特征空间，记 \(p_{\lambda}\) 为像为 \(\mathrm{E}_{\lambda}\) 的正交投影；它是算子 \(u\) 关于 \(\{\lambda\}\) 的谱投影（IV 第 278 页命题 9 的推论）。根据敏感谱的定义，投影 \(p_{\lambda}\) 具有有限秩，因而是紧的。置 \(w = u + p_{\lambda}\)；它是自伴部分算子。为完成证明，我们验证 \(\lambda\) 属于 \(w\) 的预解集。

有 \(\mathrm{E}_{\lambda} \subset \operatorname{dom}(u)\)，且空间 \(\mathrm{E}_{\lambda}\) 和 \(\operatorname{dom}(u) \cap \mathrm{E}_{\lambda}^{\circ}\) 在 \(u\) 下稳定（IV 第 278 页命题 9）。

设 \(x \in \text{dom}(u)\)。写成 \(x = p_{\lambda}(x) + y\)，其中 \(y \in \text{dom}(u) \cap \text{E}_{\lambda}^{\circ}\)。由上文可得

\[
\left\| w (x) - \lambda x \right\| ^ {2} = \left\| p _ {\lambda} (x) \right\| ^ {2} + \left\| w (y) - \lambda y \right\| ^ {2}.\tag{21}
\]

令 V 为不与 u 的谱相交的 \(\lambda\) 的一个开邻域，并令 c > 0，使得以 \(\lambda\) 为中心、半径为 c 的圆盘包含于 V。令 \(\mu_{y}\) 为 y 关于 u 的谱测度。由于 y 与 \(E_{\lambda}\) 正交，\(\mu_{y}\) 的支撑不与 V 相交（同上）。于是有

\[
\| w (y) - \lambda y \| ^ {2} = \| u (y) - \lambda y \| ^ {2} = \int_ {\mathbf {C}} | t - \lambda | ^ {2} d \mu_ {y} (t) \geqslant c ^ {2} \| y \| ^ {2}.\tag{22}
\]

由 (21) 和 (22)，我们得到 \(\|w(x)-\lambda x\|^{2}\geqslant\inf(c^{2},1)\|x\|^{2}\)，因为 w 是自伴的且 \(\lambda\in R\)。结论于是由 IV 第 248 页引理 5 得出。

推论。 --- 设 \(u\) 是 E 上的自伴部分算子，\(v\) 是 E 的紧埃尔米特自同态。有 \(\mathrm{Sp}_{\mathrm{e}}(u + v) = \mathrm{Sp}_{\mathrm{e}}(u)\)。

这立即由定理得出。

